% Content moved out of Section 2 when the main text was brought within the Biometrics page limit.
% Included from sections/appendix.tex after Web Appendix A.8.

\subsection{A.9 What standard estimators estimate under feedback}
\setreflabel{A.9}
\label{ss:bias}

Parts (a)--(c) take $\kappa_0=\phi_0=0$ and the pressure $\Lambda_{it}=L_{it}+\omega_i$ as
known, which favors the comparators. $\ddot D_{it}$ denotes two-way demeaned $D_{it}$ in a
balanced panel, $d^2=\lim N^{-1}\E\sum_{i,t}\ddot D_{it}{}^2>0$, and limits are taken as
$G\to\infty$ with $T$ fixed.

\begin{theorem}[Pseudo-true values]
\label{thm:bias}
Assume A1--A6 and (L).
\begin{enumerate}
\item[(a)] \emph{Two-step log-$R$ difference-in-differences.} Let $\hat\tau_2$ be the two-way
fixed-effects slope of $y_{it}=\log\{(Y_{it}+c)/\Lambda_{it}\}$ on $D_{it}$, $c>0$
\citep{feng2024parallel}. Define $b_{it}=\log\mu_{it}-\E\{\log(Y_{it}+c)\mid\F_{t-1}\}$ and
$\eta_{it}=\log(Y_{it}+c)-\E\{\log(Y_{it}+c)\mid\F_{t-1}\}$. Then exactly
$\hat\tau_2-\tau_0=\{-\sum\ddot D_{it}b_{it}+\sum\ddot D_{it}\eta_{it}\}/\sum\ddot D_{it}{}^2$, and
$\hat\tau_2\to_p\tau_0+B_J+B_F$ with
\begin{gather*}
B_J=-\lim\frac{\E\sum_{i,t}\ddot D_{it}b_{it}}{Nd^2},\\
B_F=-\lim\frac{\E\sum_i\sum_{t<s}D_{is}\eta_{it}}{NTd^2},\qquad
|B_F|\le\frac{2\bar\sigma_\eta T^{1/2}}{3d^2},
\end{gather*}
where $\bar\sigma_\eta{}^2=\sup_{i,t}\E\eta_{it}{}^2$. Under a negative binomial working law with mean
$\mu$ and shape $k$, a second-order expansion gives $b_{it}\approx-(c-1/2+c/k)/\mu_{it}+1/(2k)$.
$B_F=0$ if $\E(\eta_{it}\mid D_{i1},\dots,D_{iT})=0$, and $B_F<0$ when adoption follows unusually
high case counts.
\item[(b)] \emph{Poisson fixed effects.} Let $\hat\beta_1$ be the conditional (Poisson
fixed-effects) estimator with offset $\log\Lambda_{it}$, unit and week effects, and regressors
$x_{it}$. Its estimating function is $S_i(\beta)=\sum_s\{x_{is}-\bar x^m_i(\beta)\}Y_{is}$, where
$\bar x^m_i=\sum_s m_{is}x_{is}/\sum_s m_{is}$, and
$\E S_i(\beta_0)=-\sum_s\E\{\bar x^m_i(\beta_0)\varepsilon_{is}\}$, which is non-zero in general when
$\Lambda_{i,s+1}$ or $D_{i,s+1}$ depends on $Y_{is}$. If $Y_{is}\mid\F_{s-1}$ is Poisson and
$\bar x^m_i=f(Y_{is},R_{is})$ with $R_{is}$ conditionally independent of $Y_{is}$, then
$\E(\bar x^m_i\varepsilon_{is}\mid\F_{s-1})=\mu_{is}\E\{f(Y_{is}+1,R_{is})-f(Y_{is},R_{is})\mid\F_{s-1}\}$
\citep{chen1975poisson}. For $T=2$, $S=1$, $D_{i1}=0$ and $D_{i2}$ a function of $Y_{i1}$ and
independent noise, the pseudo-true week-2 log growth $\tilde h_d$ in group $D_{i2}=d$ solves
\begin{equation}
\E\Bigl[\ind(D_{i2}=d)\,\Lambda_{i2}\,\frac{\Lambda_{i1}e^{\alpha_i+h_0+\tau_0d}-e^{\tilde h_d}Y_{i1}}{\Lambda_{i1}+e^{\tilde h_d}\Lambda_{i2}}\Bigr]=0,
\label{eq:t2fe}
\end{equation}
and $\hat\tau_1\to_p\tilde h_1-\tilde h_0$. Because $\Lambda_{i2}=Y_{i1}+\omega_i$, $\tilde h_d\ne h_0+\tau_0d$, and in general
$\tilde h_1-\tilde h_0\ne\tau_0$, even when $D_{i2}$ is randomized.
\item[(c)] \emph{Two-way fixed effects on $\log(Y+1)$.} The slope $\hat\tau_3$ satisfies
\[
\hat\tau_3-\tau_0=\frac{\sum\ddot D_{it}\log\Lambda_{it}}{\sum\ddot D_{it}{}^2}+\frac{-\sum\ddot D_{it}b^{(1)}_{it}+\sum\ddot D_{it}\eta^{(1)}_{it}}{\sum\ddot D_{it}{}^2},
\]
with $b^{(1)},\eta^{(1)}$ as in (a) with $c=1$. The first, omitted-pressure term is positive when
adoption follows rising pressure and negative when the policy lowers later pressure, so the
sign of the bias is not determined.
\item[(d)] \emph{Ignoring spillover.} Let $(\tilde\delta,\tilde\tau)$ be the pseudo-true value of
QD-GMM with $\kappa=\phi=0$ imposed and instruments $s_{it}e_t$, $s_{it}D_{it}$. It solves
$\lim G^{-1}\E\sum_{i,t}z_{it}\mu_{it}\{R_{it}(\delta,\tau)-1\}=0$, where
\begin{multline*}
R_{it}(\delta,\tau)=\exp\{\gamma_{t+1}-\gamma_t-\delta_t\\
-(\tau-\tau_0)\Delta D_{i,t+1}+\kappa_0\Delta\bar D_{i,t+1}\}\,
\frac{1+\phi_0\tilde N_{i,t+1}}{1+\phi_0\tilde N_{it}},
\end{multline*}
$\tilde N_{it}=N_{it}/(L_{it}+\omega_i)$ and $\Delta X_{i,t+1}=X_{i,t+1}-X_{it}$. If $\kappa_0=\phi_0=0$ then $\tilde\tau=\tau_0$. To first
order in $(\kappa_0,\phi_0)$, $\tilde\tau-\tau_0$ is the coefficient of $\Delta D_{i,t+1}$ in the
just-identified instrumental-variable fit, with weights $s_{it}\mu_{it}$ and instruments
$(e_t,D_{it})$, of $\kappa_0\Delta\bar D_{i,t+1}+\phi_0\Delta\tilde N_{i,t+1}$ on week indicators
and $\Delta D_{i,t+1}$. When $\Delta\bar D$ moves with $\Delta D$, as under state-level adoption,
the $\kappa_0$ contribution has the sign of $\kappa_0$.
\end{enumerate}
\end{theorem}

Theorem~\ref{thm:bias} adapts known results on count panels with predetermined regressors to the
renewal mean; its contribution is to place the comparators on a common footing, not to establish new
inconsistency results. Part (a) separates two mechanisms. The Jensen term $B_J$ is the small-count problem of
estimated reproduction numbers and does not vanish as $T$ grows. The feedback term $B_F$ arises
because the unit mean of $D$ uses future adoption, which responds to current shocks. This is the
mechanism behind the fixed-$T$ bias of within estimators in dynamic panels
\citep{nickell1981biases}. Because $d^2=\lim N^{-1}\E\sum_{i,t}\ddot D_{it}{}^2$ sums across all
$t=1,\dots,T$, in designs with non-vanishing within-unit treatment variation per period $d^2$ scales
linearly with $T$. The general bound in (a) is therefore $O(T^{-1/2})$, and under the recent-shocks
condition with cross-product sum $O(T)$, the bound $C/d^2$ is $O(T^{-1})$, matching the Nickell
rate.
Part (b) illustrates the classical fixed-$T$ inconsistency of the Poisson fixed-effects estimator
\citep{hausman1984count} when regressors, here the lagged-outcome offset, are predetermined rather
than strictly exogenous \citep{wooldridge1997multiplicative,wooldridge1999distribution,blundell2002count};
it holds even with randomized adoption, but in the randomized designs of Web Appendix~\ref{app:numeric}
the bias is $-0.004$, so (b) is an illustration rather than a practical concern there. Part (c) is
the log-linearization problem of \citet{santossilva2006log} and \citet{chen2024logs} combined with
an omitted pressure term; Poisson-based difference-in-differences \citep{wooldridge2023nonlinear}
removes the former but not the latter.
Split-panel jackknife corrections \citep{dhaene2015splitpanel} assume stationarity, which week
effects and staggered adoption violate, and analytical corrections for nonlinear panels with large
$N$ and $T$ \citep{fernandezval2016individual} require long panels, so the corrected estimator is
assessed only by simulation.

\subsection{A.10 Heterogeneous and dynamic effects}
\setreflabel{A.10}
\label{ss:hetero}

\begin{proposition}[Heterogeneous effects]
\label{prop:hetero}
Suppose A1 holds with unit-specific, time-invariant $\tau_i$, possibly dependent on $\alpha_i$ and
on the adoption process, adoption is absorbing, and $(\kappa,\phi)=(\kappa_0,\phi_0)$ is fixed. Use
the instruments of Theorem~\ref{thm:ident}(b). Let $M_t=\lim G^{-1}\E\sum_i s_{it}\mu_{it}$ and
assume $A_t>0$, $K_t>0$ for some $t$. Then the pseudo-true value $\tilde\tau$ exists, is unique,
and lies in the closed convex hull of the support of $\tau_i$ under the measure
$s_{it}\mu_{it}\Sw_{i,t+1}$ restricted to weeks with $A_t>0$. A first-order expansion in
$\tau_i-\tilde\tau$ gives the following approximation.
\[
\tilde\tau\approx\frac{\sum_t(A_t/M_t)\lim G^{-1}\E\sum_i s_{it}\mu_{it}\Sw_{i,t+1}\tau_i}{\sum_t(A_t/M_t)K_t},
\]
This is a weighted average with non-negative weights.
\end{proposition}

Proposition~\ref{prop:hetero} adapts the weighting analyses of two-way fixed effects with
staggered adoption \citep{goodmanbacon2021did,dechaisemartin2020twfe} to the quasi-differenced
moment; it holds $(\kappa,\phi)$ at their true values and its weights are first order. Unlike two-way
fixed effects, the constant-effect QD-GMM estimand cannot place negative weight on any unit's effect
when effects are time-invariant. Switchers in weeks
with no treated units receive zero first-order weight.

\begin{remark}[Dynamic effects]
\label{rem:dynamic}
If the effect changes with time since adoption, treated stayers no longer cancel $\tau$; their
case growth contains $\tau_{e+1}-\tau_e$. Because already-treated units help pin down the week
effects (Proposition~\ref{prop:sources}), the constant-effect estimand can lie outside the range of
the true effects and can even have the wrong sign, as with already-treated comparisons in
two-way fixed effects \citep{goodmanbacon2021did}. We therefore recommend the event-time
specification $\sum_{e=0}^{E}\tau_e\ind\{t-A_i=e\}$, where $A_i$ is the adoption week and the
last bin collects $e\ge E$,
instrumented by $s_{it}\ind\{t-A_i=e\}$. These indicators are $\F_{t-1}$-measurable, so
Theorem~\ref{thm:asy} applies whenever the rank condition holds. The rank condition requires
not-yet-treated or never-treated units.
\end{remark}

Remark~\ref{rem:dynamic} adapts the event-study analysis of \citet{sun2021estimating} and
\citet{borusyak2024imputation} to the quasi-differenced moments. The implicit weights can be made
explicit. Suppose the event-time model holds with bin effects $\tau^*=(\tau^*_0,\dots,\tau^*_E)^{\top}$
and let $\bar g_P(\boldsymbol\tau,\kappa,\phi)$ denote the limit of the pooled just-identified moments
(week effects profiled out) when the policy enters through bin-specific coefficients
$\boldsymbol\tau$; the pooled specification sets $\boldsymbol\tau=\tau\mathbf 1$. The pooled
pseudo-true value $(\tilde\tau,\tilde\kappa,\tilde\phi)$ solves $\bar g_P(\tilde\tau\mathbf
1,\tilde\kappa,\tilde\phi)=0$. With $H_b=\partial\bar g_P/\partial\boldsymbol\tau^{\top}$,
$H_\kappa=\partial\bar g_P/\partial\kappa$ and $H_\phi=\partial\bar g_P/\partial\phi$ at
$(\tau^*,\kappa_0,\phi_0)$, a first-order expansion gives
\[
A\,(\tilde\tau,\ \tilde\kappa-\kappa_0,\ \tilde\phi-\phi_0)^{\top}=H_b\tau^*,\qquad
A=[H_b\mathbf 1,\ H_\kappa,\ H_\phi],
\]
so that $\tilde\tau\approx\sum_{b}\omega_b\tau^*_b$, where $\omega^{\top}$ is the first row of
$A^{-1}H_b$. Because $A^{-1}H_b\mathbf 1=e_1$, the weights sum to one; they are not restricted to be
non-negative, because treated stayers enter $H_b$ with the opposite sign to switchers. The second row
shows that $\tilde\kappa$ is also contaminated when the bin effects differ. Web
Table~\ref{tab:rev-event} reports the dynamic event-time profile around summer holiday onset for
the application of Section~\ref{s:application}.

\subsection{A.11 Cases averted}

Set $\nu_{it}=Y_{it}$ for $t\le 0$. For $t\ge 1$ the counterfactual is defined recursively, for all units jointly, as follows.
\begin{equation}
\nu_{it}(\theta)=c_{it}(\theta)Y_{it},\qquad
c_{it}(\theta)=e^{-\tau D_{it}-\kappa\bar D_{it}}\,\frac{\Lambda_{it}(\nu;\phi)}{\Lambda_{it}(Y;\phi)},
\label{eq:nu}
\end{equation}
Here $\Lambda_{it}(\nu;\phi)$ is the pressure computed from $\nu$. The recursion leaves
$\nu_{it}=Y_{it}$ until a policy affects unit $i$ directly, through $\bar D_{it}$, or through the
pressure it receives, so no counterfactual start date has to be chosen. Let $\hat\psi_{\mathcal P}=\sum_{(i,t)\in\mathcal P}\{\nu_{it}(\hat\theta)-Y_{it}\}$
and $\psi^\dagger_{\mathcal P}=\sum_{(i,t)\in\mathcal P}\{\nu_{it}(\theta_0)-Y_{it}\}$.

\begin{theorem}[Counterfactual propagation]
\label{thm:psi}
(a) Under A1--A3 and A5, $\E(\nu_{it}(\theta_0)\mid\alpha,Y_{\le 0})=\E\{Y_{it}(\mathbf 0)\mid\alpha,Y_{\le 0}\}$
for all $i$ and $t$, so $\E\psi^\dagger_{\mathcal P}=\psi_{\mathcal P}$.
(b) Under the conditions of Theorem~\ref{thm:asy}, if $|\mathcal P|/G$ is bounded and the
dominance conditions of Web Appendix A.1 hold for $\partial\nu_{it}/\partial\beta$, then
$\sqrt G(\hat\psi_{\mathcal P}-\psi^\dagger_{\mathcal P})/G\to_d N(0,\Psi V\Psi^{\top})$ with
$\Psi=\plim G^{-1}\sum_{(i,t)\in\mathcal P}\partial\nu_{it}(\theta_0)/\partial\theta^{\top}$.
\end{theorem}

The naive rescaling $Y_{it}e^{-\tau D_{it}-\kappa\bar D_{it}}$ ignores that averted cases would
themselves have infected others; Theorem~\ref{thm:psi} propagates the counterfactual through the
renewal equation and the network. Part (b) is a standard delta-method consequence of
Theorem~\ref{thm:asy} and adds no new asymptotic argument. Inference in (b) targets the realized-path quantity
$\psi^\dagger$, whose expectation is $\psi$.

\begin{remark}[Process noise and depletion]
\label{rem:psi}
$\psi^\dagger_{\mathcal P}-\psi_{\mathcal P}$ is a sum of mean-zero, locally dependent terms.
Unlike the moment contributions, these terms are correlated across neighboring clusters, because
$\nu_{it}$ depends on neighbors' outcomes. A variance estimator for the population-level
$\psi_{\mathcal P}$ would add a cluster-graph HAC term; we conjecture that it is consistent but do
not prove it. If susceptible depletion multiplies the mean in A1 by
$1-C_{i,t-1}/P_i$, where $C$ is cumulative cases and $P$ population, the recursion is no longer
linear in past $\nu$ and Theorem~\ref{thm:psi}(a) holds only approximately; depletion also makes
effect estimates drift with time since exposure \citep{janvin2025depletion}.
\end{remark}

\begin{remark}[Testability]
\label{rem:testable}
Lemma~\ref{lem:mds} implies that moment contributions of adjacent clusters are uncorrelated,
which can be checked directly. In the application, the correlations of the state score
contributions over the 29 adjacent German state pairs are 0.175 ($\tau$), 0.123 ($\kappa$) and $-0.119$
($\phi$), with $-0.073$ on the total-policy influence score $\tau+\kappa$. Formal standardized
pair-product statistics ($Z_{A4}=0.966$, $0.903$, $-0.520$, and $-0.295$; $B=2000$ randomization
$p=0.350$, $0.403$, $0.627$, and $0.767$) are indistinguishable from zero noise at $G=16$. Adding the
unshrunk adjacent-pair cross-products to $\hat\Omega$ produces a non-positive-semidefinite
adjusted covariance matrix ($\min\mathrm{eig}=-0.016$ on $\hat V_b$, with a negative diagonal
variance entry on $\tau$), even as the formal quadratic form for $\hat\tau+\hat\kappa$ drops from
0.0584 to 0.0500 and for $\hat\phi$ from 0.0566 to 0.0356; under Assumption~A4$'$, eigenvalue-clipped
PSD spatial adjustment widens SE($\hat\tau+\hat\kappa$) from $0.0584$ to $0.0703$ and
macro-region ($G=4$) clustering gives $0.0617$. Because the unshrunk sample HAC matrix is indefinite,
it provides no basis for treating the ordinary state-clustered sandwich as conservative (consistent
with the calibrated simulation in Web Table~\ref{tab:rev-calsim}, where sandwich SEs understate
empirical SD by 15.0\%, SE ratio $0.850$, driven by cluster-size skewness $\times T$ in Web
Table~\ref{tab:rev-clustergrid}). The
$J$ test of Corollary~\ref{cor:J}, in the spirit of \citet{arellano1991tests}, checks the
multiplicative renewal structure against extra predetermined instruments. A sensitivity analysis
lagging $D$ by one week checks sensitivity to policy-timing alignment. Sequential ignorability
(A2(b)) and parallel untreated log-transmission paths after adoption are not testable; under
feedback, lead regressions are not valid tests (Proposition~\ref{prop:sources}(c)), pre-testing
distorts subsequent inference \citep{roth2022pretest}, and bounded departures from parallel trends
\citep{rambachan2023credible} or a design-based view of dynamic effects \citep{bojinov2021panel}
offer alternatives that we do not pursue.
\end{remark}

\subsection{A.12 Transmission within a week and assumption A4}
\setreflabel{A.12}
\label{ss:weekwithin}

A4 states that, conditionally on the history, outcomes are independent across clusters within a
week. If an infection crosses a cluster boundary inside the same week and both counties are
recorded in that week, the two week-$t$ innovations are dependent given $\F_{t-1}$, and A4 fails
exactly where it matters. When does that happen, and what does it cost?

First, the assumption is about the innovations, not the outcomes. Neighboring clusters share
history through $\bar D$ and through the network pressure $\Lambda$, and Lemma~\ref{lem:mds}
accounts for that by weighting each moment contribution with a $\F_{s-1}$-measurable factor. What
A4 excludes is a correlation that arrives in the same week as the outcome, before it can enter
the information set.

Second, a weekly panel can only be internally consistent if contagion is lagged at least one
period, because the renewal equation used to define the estimand, equation~\eqref{eq:mean}, is
written in completed weeks. Same-week transmission therefore changes the conditional mean implied
by that model and may also change the covariance structure of the quasi-difference. The resulting
moments need not be martingale differences, so neither the consistency argument nor the conventional
cluster-robust sandwich follows without additional conditions. The direction and magnitude of parameter and variance distortions are not determined by the weekly
model alone. For fast-generation pathogens such as SARS-CoV-2, roughly 30\% of transmission falls
within the same week; the graded grid and matched-control experiments in Section~\ref{ss:sameweek}
generate within-week own-unit renewal across $w_0\in[0,0.30]$ and aggregate to weeks. At $w_0=0.30$
(14.2\% nonlinear within-week reinfection share, or 34.0\% including seeded same-week incidence),
bias on $\tau+\kappa$ reaches $-0.053$ (coverage 0.783; 0.890 for $\tau$) in the unadjusted arm and
$-0.046$ (coverage 0.859) under treated-share matching, with an unpowered $J$ test (4.8\%). By
contrast, for longer-generation pathogens such as varicella in Section~\ref{s:application}, the
generation interval places less than $10^{-5}$ of its mass within the same week
($w_0=2.6\times 10^{-6}$), so the within-week aggregation bias mechanism is negligible.

Third, dependence created by same-week transmission can link innovations across clusters before
the next period's information set is observed. This is distinct from the lagged cross-cluster
transmission allowed by A4 and Lemma~\ref{lem:mds}; the existing conditions do not establish that
the resulting moment contributions remain uncorrelated beyond the instrument radius. A spatial
correction such as that of \citet{conley1999gmm} addresses covariance under appropriate dependence
conditions, but it does not by itself repair a misspecified conditional mean or restore the
martingale-difference property. The 25 pre-specified scenarios do not include same-week transmission; the separate matched
simulation of Section~\ref{ss:sameweek} quantifies its effect on bias and coverage.

Fourth, weekly aggregation may be inadequate when a substantial share of transmission occurs within
a week. Shortening the observation interval or specifying a within-week renewal process would change
the data model: daily counts carry day-of-week reporting effects and reporting delays
\citep{lison2024nowcast}, which violate the delay-free reporting of
Proposition~\ref{prop:ascertainment} and require new moment conditions and new simulation
validation. Variance corrections alone cannot resolve the mean specification, so fast-generation pathogens with
large same-week mass ($w_0\gtrsim 0.3$) require a dedicated sub-weekly estimator.

In the same-week simulation, the within-week mechanism is switched on from the third window week,
that is, from index 2 counting the first window week as 0.

